\begin{exo}{}
	Le prix moyen des appartements anciens dans $7$ villes françaises est donné dans le tableau ci-dessous.
	\begin{center}
		\footnotesize
		\begin{tabular}{|c|c|c|c|c|c|}\hline
			                     & \cellcolor{gray!30}A & \cellcolor{gray!30}B & \cellcolor{gray!30}C & \cellcolor{gray!30}D                                                 & \cellcolor{gray!30}E                                                  \\ \hline
			\cellcolor{gray!30}1 & \textbf{Ville}       & \textbf{2016}        & \textbf{2017}        & \textbf{\begin{tabular}[c]{@{}c@{}}variation\\ absolue\end{tabular}} & \textbf{\begin{tabular}[c]{@{}c@{}}variation\\ relative\end{tabular}} \\ \hline
			\cellcolor{gray!30}2 & Strasbourg           & $\np{3279}$          & $\np{3531}$          &                                                                      &                                                                       \\ \hline
			\cellcolor{gray!30}3 & Bordeaux             & $\np{4155}$          & $\np{4599}$          &                                                                      &                                                                       \\ \hline
			\cellcolor{gray!30}4 & Caen                 & $\np{2307}$          & $\np{2390}$          &                                                                      &                                                                       \\ \hline
			\cellcolor{gray!30}5 & Paris                & $\np{8967}$          & $\np{9559}$          &                                                                      &                                                                       \\ \hline
			\cellcolor{gray!30}6 & Nîmes                & $\np{2124}$          & $\np{2381}$          &                                                                      &                                                                       \\ \hline
			\cellcolor{gray!30}7 & Orléans              & $\np{2158}$          & $\np{2111}$          &                                                                      &                                                                       \\ \hline
			\cellcolor{gray!30}8 & Lyon                 & $\np{4279}$          & $\np{4604}$          &                                                                      &                                                                       \\ \hline
		\end{tabular}
	\end{center}

	Les cellules de la colonne \textbf{E} sont au format pourcentage écrit avec une décimale (\textit{source : baromètre LPI-SeLoger}).
	\begin{enumerate}
		\item Quelles formules peut-on entrer dans les cellules \textbf{D3} et \textbf{E3} pour que, une fois recopiées vers le bas, elles affichent les variations absolue et relative pour chacune des villes ?
		\item Reproduire et compléter le tableau.
		\item Dans laquelle de ces villes le prix connaît-il la plus forte augmentation en euros ? en pourcentage ?
	\end{enumerate}
\end{exo}
